\documentclass[11pt,a4paper]{article}

\usepackage{../../shared/cathedral-preamble}
\usepackage{setspace}
\onehalfspacing

\title{%
  The Cathedral: A Guide for the Curious\\[4pt]
  \large What a Computer Proof About Prime Numbers\\
  \large Means for the Rest of Us%
}
\author{Jason Robert Gochanour}
\date{June 2026}

\begin{document}
\maketitle

\begin{abstract}
This paper explains the Cathedral project---a computer-verified
proof about the Riemann Hypothesis---for readers without a
mathematics background. No equations are required. We explain what
the Riemann Hypothesis is, why it matters, what the Cathedral
actually proved, what we accidentally discovered about the
connection between prime numbers and physics, and why the answer
is ``we reduced a 167-year-old mystery to two well-understood
facts, a computer checked our work, and then the integers
showed us they have a physics.''
\end{abstract}

% ================================================================
\section{The Oldest Open Problem}
% ================================================================

There is a question about prime numbers---2, 3, 5, 7, 11, 13,
\ldots---that has been open since 1859. A German mathematician
named Bernhard Riemann conjectured a pattern in how primes are
distributed among the whole numbers. His conjecture, called the
\emph{Riemann Hypothesis}, has never been proved or disproved.

It is considered the most important unsolved problem in mathematics.
A proof would have implications for cryptography, physics, and our
understanding of the fundamental structure of numbers.

% ================================================================
\section{What We Did}
% ================================================================

We did \textbf{not} prove the Riemann Hypothesis.

What we did is this: we built a formal proof---checked by a
computer, line by line---that the Riemann Hypothesis is
\emph{exactly equivalent} to a different mathematical statement.
This other statement says that a certain distance measurement
(called the Nyman--Beurling distance) gets closer and closer to
zero.

Think of it like this: imagine you want to know whether a mountain
has gold inside it. You can't dig through the mountain. But you
\emph{can} prove, with mathematical certainty, that the mountain
has gold \emph{if and only if} a certain needle on a gauge points
to zero. You haven't found the gold, but you've built a perfect
measuring instrument.

Then, while building the instrument, you discovered that the
mountain has the same internal structure as a particle
accelerator. That discovery is almost as interesting as the
gold.

% ================================================================
\section{What Does ``Computer-Verified'' Mean?}
% ================================================================

When a human writes a mathematical proof, other humans check it.
Humans make mistakes. In 2012, a proof of a major conjecture was
published, celebrated, and later found to contain an error.

When a computer verifies a proof, it checks every logical step
against the rules of mathematics, mechanically, without skipping
anything. It cannot be fooled, bribed, or tired. If the computer
says the proof is correct, it is correct---up to the correctness
of the computer itself.

We used a program called \textbf{Lean~4}, developed at Microsoft
Research, which has been used to verify proofs in algebra, geometry,
and analysis. Our proof spans $\sim$120{,}000 lines across 390 active
files and compiles without errors.

% ================================================================
\section{What Remains}
% ================================================================

Our proof reduces the Riemann Hypothesis to \textbf{two} remaining
facts:

\begin{enumerate}
  \item \textbf{The Gram form bound}: The Riemann Hypothesis itself,
    reformulated as a statement about a specific table of numbers
    (the Gram matrix). This is the open question.
  \item \textbf{The Mertens bound}: A known result about how well
    the M\"obius function cancels. This has been proved by humans
    since 1896 and is unconditionally true---it is only an ``axiom''
    because the computer's math library doesn't yet include it.
\end{enumerate}

Everything else---3{,}500+ theorems across 390 files---has been
verified by the computer with zero human trust required.

% ================================================================
\section{The Surprise: Primes Have a Physics}
% ================================================================

While building the proof, we accidentally discovered that the
prime numbers behave like particles in a quantum field.

This is not a vague analogy. The proof structures correspond
precisely to structures in physics:

\begin{itemize}
  \item The table of numbers at the heart of the proof (the Gram
    matrix) behaves exactly like a quantum correlator---the tool
    physicists use to measure how particles interact.
  \item The prime numbers play the role of independent dimensions
    in an infinite-dimensional geometric space, like the hidden
    dimensions in string theory.
  \item The proof's convergence mechanism---the reason the distance
    measurement shrinks toward zero---works exactly like the
    ``asymptotic freedom'' that governs quarks inside protons.
  \item The alternating signs of the M\"obius function (a classical
    number theory tool) act like Pauli exclusion in quantum
    mechanics---the principle that keeps matter from collapsing.
\end{itemize}

These are not metaphors. They are machine-verified theorems.
A 63-page companion paper documents over 160 such correspondences.

Why do prime numbers ``know'' physics? We don't know. But the
Cathedral makes the question precise, and that is new.

% ================================================================
\section{The Three Builders}
% ================================================================

The Cathedral was built by three minds working together:

\begin{enumerate}
  \item A human mathematician (\textbf{Jason Robert Gochanour}),
    who provided the vision, selected the approach, and made the
    judgment calls that no AI could make.
  \item Google's \textbf{Gemini} AI, which suggested mathematical
    strategies, reviewed papers, and identified the deep
    connections to physics.
  \item Anthropic's \textbf{Claude} AI, which wrote the
    computer-verified proofs, built the numerical experiments,
    and produced all $\sim$120{,}000 lines of code.
\end{enumerate}

This is, to our knowledge, the first time two AI systems and a
human have collaborated to produce a formally verified mathematical
result of this scale. The project took 68 days.

No single participant---human or AI---could have built the Cathedral
alone. The human provided intuition and taste. The AIs provided
speed and precision. The compiler provided incorruptible verification.

% ================================================================
\section{Why Does This Matter?}
% ================================================================

\begin{enumerate}
  \item \textbf{For mathematics}: It shows that computer verification
    can tackle the hardest open problems, not just textbook exercises.
    The logical structure of the Riemann Hypothesis is now mapped.

  \item \textbf{For physics}: The discovery that prime numbers mirror
    quantum field theory opens a new frontier. If this structural
    correspondence has a deep explanation, it could illuminate both
    number theory and fundamental physics.

  \item \textbf{For engineering}: The Cathedral produced four
    universal design patterns for building reliable systems---from
    sensor networks to drug discovery to parallel computing.

  \item \textbf{For AI}: It demonstrates a model where AI augments
    human capability rather than replacing it. Three minds, each
    contributing what it does best.

  \item \textbf{For security}: The same mathematics that can protect
    a power grid can, in adversarial hands, attack one. We wrote a
    responsible disclosure paper (``The Dark Mirror'') identifying
    seven dual-use risks and recommending mitigations---before
    anyone else discovers them.

  \item \textbf{For everyone}: It means the logical structure of one
    of humanity's deepest mathematical mysteries is now mapped,
    verified, and available for inspection. The assumptions are
    explicit. The reasoning is machine-checked. The trust
    boundaries are visible.
\end{enumerate}

% ================================================================
\section{For Journalists}
% ================================================================

\subsection{One-Sentence Summary}

A human-AI team built a 120{,}000-line computer-verified proof
that reduces the 167-year-old Riemann Hypothesis to two
well-understood mathematical facts---and accidentally discovered
that prime numbers have the same mathematical structure as
quantum physics.

\subsection{Key Facts}

\begin{itemize}
  \item \textbf{What}: Formal proof that RH is equivalent to a
    distance measurement going to zero, plus a structural
    correspondence between prime numbers and quantum field theory.
  \item \textbf{Who}: Jason Robert Gochanour with Google's Gemini
    and Anthropic's Claude.
  \item \textbf{How long}: 68 days, March--June 2026.
  \item \textbf{Size}: $\sim$120{,}000 lines of verified proof code
    (390 active files), 55+ computational experiments, 17 papers.
  \item \textbf{What it is NOT}: A proof of the Riemann Hypothesis.
  \item \textbf{What it IS}: A machine-verified reduction to two
    standard results, plus a physics dictionary with 160+
    verified correspondences.
\end{itemize}

\subsection{Suggested Headlines}

\emph{``Computer Proof Maps the Logical Structure of Math's Greatest
Unsolved Problem---and Discovers Prime Numbers Have a Physics''}

\emph{``Human-AI Team Reduces 167-Year-Old Math Mystery to Two
Facts, Finds Primes Behave Like Quantum Particles''}

\subsection{Common Misconceptions}

\begin{enumerate}
  \item \textbf{``They proved RH!''} --- No. We proved an equivalence
    and reduced it to two well-understood facts.
  \item \textbf{``AI solved a math problem.''} --- AI assisted; the
    human directed and verified. Neither could have done it alone.
  \item \textbf{``This is just a computer program.''} --- It is a
    proof, checked by a theorem-proving compiler. If the compiler
    accepts it, the mathematics is correct.
  \item \textbf{``The physics connection is just a metaphor.''} ---
    No. It consists of 160+ machine-verified structural
    correspondences documented in a 63-page paper.
\end{enumerate}

% ================================================================
\section{How to Verify}
% ================================================================

The entire proof is open source. To check it yourself:

\begin{enumerate}
  \item Install Lean~4 (\url{https://leanprover.github.io}).
  \item Clone the repository.
  \item Run \texttt{cd proofs \&\& lake build Cathedral}.
  \item If it compiles: the proof is correct.
\end{enumerate}

No trust in the authors is required. The compiler is the arbiter.

\begin{thebibliography}{99}
\bibitem{cathedral}
J.R.~Gochanour, \emph{The Cathedral}, 2026.
\bibitem{cathedral-physics}
J.R.~Gochanour, \emph{The Physics of the Primes}, 2026.
\end{thebibliography}

\end{document}
